\section{Extensiones finitas y algebraicas; fórmula de la torre}\label{sec:campos-algebraicas}

\subsection{Grados, bases y fórmula de la torre}\label{tema:23-1-1}\label{tema:23-1-2}
\begin{defi}
Sea $k\subseteq F$ una extensión de campos. Un elemento $\alpha\in F$ es \emph{algebraico sobre $k$} si $[k(\alpha):k]$ es finito; ese entero es su \emph{grado sobre $k$}. Si el grado es infinito, $\alpha$ es \emph{trascendente sobre $k$}. La extensión es \emph{algebraica} si todos sus elementos son algebraicos sobre $k$.
\end{defi}
Por la dicotomía de la evaluación de la \cref{prop:C22-simple}, $\alpha$ es algebraico exactamente cuando existe $0\ne f\in k[t]$ con $f(\alpha)=0$. Su polinomio mínimo es el único mónico de menor grado que lo anula; es irreducible y su grado es $[k(\alpha):k]$.

\begin{lema}\label{lema:C23-finita-algebraica}
Si $[F:k]=n<\infty$, todo elemento de $F$ es algebraico sobre $k$, de grado a lo sumo $n$.
\end{lema}
\begin{proof}
Para $\alpha\in F$, el subcampo $k(\alpha)$ es un subespacio vectorial de $F$ sobre $k$, luego su dimensión es a lo sumo $n$. Concretamente, los $n+1$ elementos $1,\alpha,\ldots,\alpha^n$ son dependientes. Una relación no trivial entre ellos proporciona un polinomio no nulo que se anula en $\alpha$.
\end{proof}

\begin{prop}[Fórmula de la torre]\label{prop:C23-torre}
Sean $k\subseteq E\subseteq F$ extensiones de campos. El grado $[F:k]$ es finito si y sólo si $[E:k]$ y $[F:E]$ lo son. En ese caso,
\[
[F:k]=[F:E]\,[E:k].
\]
Si $(e_1,\ldots,e_m)$ es una base de $E$ sobre $k$ y $(f_1,\ldots,f_n)$ una base de $F$ sobre $E$, los $mn$ productos $e_if_j$ forman una base de $F$ sobre $k$.
\end{prop}
\begin{proof}
Si $[F:k]=N$ es finito, el subespacio $E$ tiene dimensión finita sobre $k$. Además, cualquier familia de $N+1$ elementos de $F$ satisface una relación con coeficientes en $k$ no todos nulos. Esa misma relación vale sobre $E$, de modo que $[F:E]\leq N$.

Recíprocamente, supongamos finitos ambos grados y tomemos las bases del enunciado. Para $g\in F$ existen $d_j\in E$ tales que $g=\sum_jd_jf_j$. Escribiendo $d_j=\sum_ic_{ij}e_i$, con $c_{ij}\in k$, obtenemos
\[
g=\sum_{j=1}^n\sum_{i=1}^m c_{ij}e_if_j.
\]
Los productos generan $F$ sobre $k$. Para probar su independencia, si
\[
0=\sum_{i,j}\lambda_{ij}e_if_j
=\sum_j\left(\sum_i\lambda_{ij}e_i\right)f_j,
\qquad \lambda_{ij}\in k,
\]
la independencia de los $f_j$ sobre $E$ da $\sum_i\lambda_{ij}e_i=0$ para cada $j$. La independencia de los $e_i$ sobre $k$ implica entonces $\lambda_{ij}=0$ para todos los índices. Por tanto la dimensión es $mn$.
\end{proof}

\begin{coro}\label{coro:C23-divisibilidad}
Si $[F:k]$ es finito y $k\subseteq E\subseteq F$, tanto $[E:k]$ como $[F:E]$ dividen a $[F:k]$.
\end{coro}
\begin{proof}
La fórmula de la torre expresa $[F:k]$ como el producto de esos dos enteros positivos.
\end{proof}
Una extensión tiene grado uno si y sólo si ambos campos coinciden: si $[F:k]=1$, el elemento $1$ es una base y todo $f\in F$ es un múltiplo de $1$ por un escalar de $k$.

\begin{ejem}[Un elemento de grado impar]
Si $\alpha$ es algebraico de grado impar sobre $k$, entonces $k(\alpha)=k(\alpha^2)$ y $\alpha$ es un polinomio en $\alpha^2$ con coeficientes en $k$.

En efecto, en la torre
\[
k\subseteq k(\alpha^2)\subseteq k(\alpha)
\]
el grado $d=[k(\alpha):k(\alpha^2)]$ es a lo sumo dos, pues $\alpha$ es raíz de $t^2-\alpha^2\in k(\alpha^2)[t]$. También divide al entero impar $[k(\alpha):k]$, luego $d=1$. Así $\alpha\in k(\alpha^2)$. Como $\alpha^2$ es algebraico, la evaluación sobre $k$ identifica $k[\alpha^2]$ con el campo $k(\alpha^2)$; por eso basta una expresión polinómica, sin denominadores.
\end{ejem}

\subsection{Generación como campo y como álgebra}\label{tema:23-2-1}
\begin{defi}
La extensión $k\subseteq F$ es \emph{finitamente generada} si existen $\alpha_1,\ldots,\alpha_r\in F$ tales que
\[
F=k(\alpha_1)(\alpha_2)\cdots(\alpha_r).
\]
Denotamos este campo por $k(\alpha_1,\ldots,\alpha_r)$.
\end{defi}
La generación como campo permite suma, producto y división por elementos no nulos. La generación como $k$-álgebra sólo utiliza suma y producto: su resultado es $k[\alpha_1,\ldots,\alpha_r]$. Tampoco debe confundirse con ser finito como $k$-módulo, que exige un conjunto finito de generadores lineales.

Si $\alpha$ es trascendente, la evaluación $\epsilon:k[t]\to k(\alpha)$ es inyectiva y su imagen es $k[\alpha]$, pero no es sobreyectiva. Por ejemplo, $\alpha^{-1}$ no es una expresión polinómica en $\alpha$, pues una igualdad $\alpha^{-1}=h(\alpha)$ daría $(th(t)-1)(\alpha)=0$. El diagrama correspondiente es
\[
\begin{tikzcd}[column sep=large,row sep=large]
k(t)\arrow[r,"\sim","\widetilde\epsilon"']&k(\alpha)\\
k[t]\arrow[u,hook]\arrow[ur,hook,"\epsilon"']&
\end{tikzcd}
\]
donde $k(t)$ permite fracciones de polinomios.


\begin{ejer}[Expresiones racionales en generadores]\label{ejer:C23-generadores}
Sea $k\subseteq F$ y sean $\alpha_1,\ldots,\alpha_r\in F$. Probar que
\[
k(\alpha_1,\ldots,\alpha_r)
=\set*{\frac{f(\boldsymbol\alpha)}{g(\boldsymbol\alpha)}\given
f,g\in k[t_1,\ldots,t_r],\ g(\boldsymbol\alpha)\ne0},
\]
donde $\boldsymbol\alpha=(\alpha_1,\ldots,\alpha_r)$. Deducir que el orden de los generadores es irrelevante.
\end{ejer}

\begin{sol}
Sea $R$ el conjunto del lado derecho. Todo campo que contiene a $k$ y a los $\alpha_i$ contiene las evaluaciones de $f,g$ y el inverso de $g(\boldsymbol\alpha)\ne0$. Por tanto $R\subseteq k(\alpha_1,\ldots,\alpha_r)$.

A la inversa, $R$ contiene a $k$ mediante $c=c/1$ y a cada $\alpha_i$ mediante $t_i(\boldsymbol\alpha)/1$. Si abreviamos las evaluaciones, sus operaciones son
\[
\frac{f_1}{g_1}+\frac{f_2}{g_2}
=\frac{f_1g_2+g_1f_2}{g_1g_2},
\qquad
\frac{f_1}{g_1}\frac{f_2}{g_2}=\frac{f_1f_2}{g_1g_2}.
\]
Los denominadores siguen siendo no nulos porque estamos en un campo. Los opuestos se obtienen cambiando el signo del numerador y, si $f_1(\boldsymbol\alpha)\ne0$, el inverso de $f_1/g_1$ es $g_1/f_1$. Así $R$ es un subcampo con los generadores requeridos y contiene al menor de ellos. La descripción es simétrica en los $\alpha_i$, por lo que su orden no importa. También muestra que es el menor subcampo que contiene a los campos $k(\alpha_i)$, llamado su \emph{compuesto}.
\end{sol}

\subsection{Generación finita y torres de extensiones algebraicas}\label{tema:23-2-2}\label{tema:23-2-3}
\begin{prop}\label{prop:C23-generada-algebraica}
Para $F=k(\alpha_1,\ldots,\alpha_r)$ son equivalentes:
\begin{enumerate}[nosep]
\item $[F:k]$ es finito;
\item $F/k$ es algebraica;
\item cada $\alpha_i$ es algebraico sobre $k$.
\end{enumerate}
Si $d_i=[k(\alpha_i):k]$, en estos casos $[F:k]\leq\prod_{i=1}^r d_i$.
\end{prop}
\begin{proof}
La primera condición implica la segunda por el \cref{lema:C23-finita-algebraica}; la segunda implica la tercera por definición. Para la implicación restante, definimos
\[
K_0=k,\qquad K_i=K_{i-1}(\alpha_i)\quad(1\leq i\leq r).
\]
El polinomio mínimo de $\alpha_i$ sobre $k$, de grado $d_i$, sigue anulándolo al considerarlo en $K_{i-1}[t]$. La división con resto por el mínimo sobre $K_{i-1}$ muestra que éste divide al primero. Por tanto $[K_i:K_{i-1}]\leq d_i$. Aplicando sucesivamente la fórmula de la torre,
\[
[F:k]=\prod_{i=1}^r[K_i:K_{i-1}]\leq\prod_{i=1}^r d_i<\infty.
\]
Para $r=0$, el campo es $k$ y el producto vacío es uno.
\end{proof}

\begin{obs}
La desigualdad puede ser estricta. Sobre $\mathbb R$, tanto $i$ como $-i$ tienen grado dos, pero $\mathbb R(i,-i)=\mathbb R(i)=\mathbb C$ tiene grado dos, menor que $2\cdot2$. La igualdad de campos se debe a que todo campo que contiene a $i$ contiene su opuesto.
\end{obs}

\begin{coro}\label{coro:C23-campo-algebraicos}
Si $k\subseteq F$ es una extensión, el conjunto
\[
E=\set{\alpha\in F\given \alpha\text{ es algebraico sobre }k}
\]
es un subcampo de $F$ que contiene a $k$.
\end{coro}
\begin{proof}
Todo $c\in k$ es raíz de $t-c$. Si $\alpha,\beta\in E$, la extensión $k(\alpha,\beta)/k$ es finita por la proposición y, por tanto, algebraica. Contiene a $\alpha\pm\beta$, $\alpha\beta$ y, si $\beta\ne0$, a $\alpha/\beta$. Todos estos elementos pertenecen a $E$, lo que prueba la estabilidad bajo las operaciones de campo.
\end{proof}
En particular, la algebraicidad de una expresión racional en elementos algebraicos se deduce sin calcular un polinomio que la anule. Por ejemplo, $\sqrt2+\sqrt3$ es algebraico sobre $\mathbb Q$ porque pertenece a la extensión finita $\mathbb Q(\sqrt2,\sqrt3)$.

\begin{ejem}[Los números algebraicos]
El conjunto
\[
\overline{\mathbb Q}=\set{z\in\mathbb C\given z\text{ es algebraico sobre }\mathbb Q}
\]
es un campo algebraico sobre $\mathbb Q$, pero no es una extensión finita. En efecto, hay polinomios irreducibles sobre $\mathbb Q$ de grados sin cota, como se demostró en la sección~\ref{tema:19-3-2}. Cada uno tiene una raíz en $\mathbb C$, por el teorema fundamental del álgebra, y esa raíz tiene el mismo grado que el polinomio. Esto contradice la cota del \cref{lema:C23-finita-algebraica} si se supone finita la extensión.

Además, $\overline{\mathbb Q}$ es numerable: $\mathbb Q[t]$ es numerable y cada polinomio no nulo tiene finitas raíces. Es infinito porque contiene a $\mathbb Q$. Como $\mathbb C$ es no numerable, los números trascendentes forman un conjunto no numerable. Probar la trascendencia de un número concreto exige otros argumentos; la de $e$ y la de $\pi$ se enuncian aquí sin demostración.
\end{ejem}

\begin{coro}\label{coro:C23-torre-algebraica}
Para $k\subseteq E\subseteq F$, la extensión $F/k$ es algebraica si y sólo si $E/k$ y $F/E$ son algebraicas.
\end{coro}
\begin{proof}
Si $F/k$ es algebraica, todo elemento de $E$ es algebraico sobre $k$. Para $\alpha\in F$, un polinomio no nulo de $k[t]$ que lo anula también pertenece a $E[t]$, de modo que $\alpha$ es algebraico sobre $E$.

Recíprocamente, sea $\alpha\in F$ algebraico sobre $E$, con polinomio mínimo
\[
t^n+e_{n-1}t^{n-1}+\cdots+e_0\in E[t].
\]
Sólo intervienen finitos coeficientes. El campo $K=k(e_0,\ldots,e_{n-1})$ es finito sobre $k$ porque cada $e_i$ es algebraico. El mismo polinomio muestra que $\alpha$ es algebraico sobre $K$, así que $K(\alpha)/K$ es finita. Por la fórmula de la torre, $K(\alpha)/k$ es finita y $\alpha$ es algebraico sobre $k$. Esto vale para cada $\alpha\in F$.
\end{proof}

\subsection{Una extensión de grado cuatro y sus automorfismos}\label{tema:23-3-0}
Sea $F=\mathbb Q(\sqrt2,\sqrt3)$ y $\theta=\sqrt2+\sqrt3$. La cota anterior da $[F:\mathbb Q]\leq4$, por lo que cinco potencias $1,\theta,\ldots,\theta^4$ son dependientes. Podemos hallar una relación concreta usando
\[
\theta^2=5+2\sqrt6,\qquad (\theta^2-5)^2=24.
\]
Así, $\theta$ es raíz de $f(t)=t^4-10t^2+1$. Sus cuatro raíces son $\pm\sqrt2\pm\sqrt3$ y
\[
f(t)=\prod_{\epsilon,\delta\in\{1,-1\}}
\bigl(t-(\epsilon\sqrt2+\delta\sqrt3)\bigr).
\]
Son distintas: una igualdad entre dos de ellas daría $\sqrt2=\pm\sqrt3$ o anularía una de esas raíces cuadradas.

\begin{ejem}[Irreducibilidad y generación simple]\label{ejem:C23-bicuadratica}
El polinomio $f$ es irreducible sobre $\mathbb Q$, el grado de $F/\mathbb Q$ es cuatro y $F=\mathbb Q(\theta)$.

Ninguna raíz de $f$ es racional, pues su cuadrado es $5\pm2\sqrt6$, que no es racional. Para descartar factores cuadráticos mónicos, consideramos las parejas de raíces. Una pareja de opuestas produce $t^2-(5\pm2\sqrt6)$, con término constante irracional. Si no son opuestas, las raíces comparten el signo de $\sqrt2$ o el de $\sqrt3$; su suma es $\pm2\sqrt2$ o $\pm2\sqrt3$, y el coeficiente lineal no es racional. No hay factores de grado uno ni dos sobre $\mathbb Q$, que son los únicos grados necesarios para una factorización propia de un polinomio de grado cuatro. Por tanto $f$ es el mínimo de $\theta$.

La torre $\mathbb Q\subseteq\mathbb Q(\theta)\subseteq F$ da
\[
4=[\mathbb Q(\theta):\mathbb Q]\leq[F:\mathbb Q]\leq4.
\]
El grado $[F:\mathbb Q(\theta)]$ es uno, luego ambos campos coinciden. En la otra torre, $\mathbb Q\subseteq\mathbb Q(\sqrt2)\subseteq F$, resulta $[F:\mathbb Q(\sqrt2)]=2$. En consecuencia, $t^2-3$ es irreducible sobre $\mathbb Q(\sqrt2)$.
\end{ejem}

\begin{ejer}[Automorfismos y raíces]\label{ejer:C23-raices-aut}
Sea $k\subseteq E$ una extensión y sea $f\in k[t]$. Probar que todo elemento de $\Aut_k(E)$ permuta el conjunto de raíces de $f$ que pertenecen a $E$.
\end{ejer}
\begin{sol}
Si $f(\beta)=0$ y $\sigma$ fija a $k$, aplicar $\sigma$ a la igualdad da $f(\sigma(\beta))=\sigma(f(\beta))=0$, pues los coeficientes quedan fijos. La aplicación inversa $\sigma^{-1}$ también preserva ese conjunto; por tanto la restricción es una permutación. Es el mismo cálculo de evaluación usado para los automorfismos de una extensión simple en el \cref{coro:C22-automorfismos}, aplicado ahora a cualquier polinomio de $k[t]$.
\end{sol}

En $F$ existe un automorfismo $s$ que fija $\sqrt2$ y cambia $\sqrt3$ por $-\sqrt3$: ambas raíces tienen el mínimo $t^2-3$ sobre $\mathbb Q(\sqrt2)$ y se aplica la extensión de isomorfismos de la \cref{prop:C22-extender-iso}. Intercambiando los papeles de $2$ y $3$, obtenemos otro automorfismo $r$ que cambia el signo de $\sqrt2$ y fija $\sqrt3$. Ambos tienen orden dos, conmutan y son distintos. Sus acciones producen las cuatro raíces de $f$:
\begingroup
\renewcommand{\arraystretch}{1.2}
\[
\begin{array}{c|cc|c}
&\sqrt2&\sqrt3&\theta\\[2pt]\hline
\id&\sqrt2&\sqrt3&\sqrt2+\sqrt3\\
r&-\sqrt2&\sqrt3&-\sqrt2+\sqrt3\\
s&\sqrt2&-\sqrt3&\sqrt2-\sqrt3\\
rs&-\sqrt2&-\sqrt3&-\sqrt2-\sqrt3
\end{array}
\]
\endgroup
Un automorfismo queda determinado por esas dos imágenes. No hay más elecciones que los dos signos independientes, luego
\[
\Aut_{\mathbb Q}(F)=\{\id,r,s,rs\}\cong C_2\times C_2.
\]
También el conteo de raíces del mínimo de $\theta$ del \cref{coro:C22-automorfismos} da exactamente cuatro automorfismos. En particular, el grupo no es cíclico.
